\section{Additive motives by presheaf localization}\label{sec:presheaf}

We construct the universal additive invariant by coefficient-linear
presheaf localization.

\subsection{Construction}

\begin{definition}\label{def:raw-presheaves}
For $A\in\nPr{(n+1)}$, define its free presheaf module $\mathsf P_n(A)$ recursively.
Put
\[
 \mathsf P_1(A)=\Fun((A^\omega)^{\op},\Sp).
\]
For $n\geq2$, form the coefficient matrix on $a,b\in A^\omega$
with entries
\begin{equation}\label{eq:raw-coefficient-hom}
 \mathcal R_A(a,b)=\mathsf P_{n-1}(\PrHom_A(a,b))
       \in\nPr{(n-1)},
 \qquad
 \mathsf P_n(A)=\Pmod_{\nPr{(n-1)}}(\mathcal R_A).
\end{equation}
Composition in $\mathcal R_A$ is induced by the lower free external product
and composition in $A$. The external products
\[
 \mathsf P_{n-1}(X)\otimes\mathsf P_{n-1}(Z)
 \longrightarrow\mathsf P_{n-1}(X\otimes Z),
 \qquad X,Z\in\nPr n,
\]
are constructed simultaneously by extending
$y_x\otimes y_z\mapsto y_{x\otimes z}$
for $x\in X^\omega$ and $z\in Z^\omega$
by colimits and coefficient actions.
Associativity and units follow inductively
from the lower free properties.

We also use this construction for a specified small coefficient
matrix $(I,B)$, where $B_{ab}\in\nPr n$ for \(a,b\in I\) and composition and units
are coefficient morphisms. At $n=1$ take spectral presheaves on the
ordinary category with mapping spaces
$\Map_{\nPr1}(\one_1,B_{ab})$.
At $n\geq2$ replace the entries in \eqref{eq:raw-coefficient-hom}
by $\mathsf P_{n-1}(B_{ab})$. Denote the result by
$\mathsf P_n(I,B)$.
\end{definition}

Informally, $\mathsf P_n(A)$ plays the role of
the category of presheaves on the underlying
$(\infty,n)$-category of compact objects of $A$,
with coefficients in $\nPr{(n-1)}$.


Restriction to represented objects gives
\begin{equation}\label{eq:raw-presheaf-up}
 \Fun^L_{\nPr{(n-1)}}(\mathsf P_n(A),E)
 \simeq\operatorname{Fun}_n^\omega(A,E),
 \qquad E\in\nPr n.
\end{equation}
The compact-preserving functors correspond to
$\operatorname{Fun}_n^{\omega,c}(A,E)$.
For $(I,B)$ the same property classifies object values indexed by
$I$ and the corresponding recursive local diagrams.

The free presheaves in Definition~\ref{def:raw-presheaves} implement
this recursion. Each local diagram extends to a compact-preserving
coefficient map
\[
 \mathsf P_{n-1}(\PrHom_A(a,b))
       \longrightarrow\PrHom_E(Fa,Fb).
\]
The compact-image requirement in \eqref{eq:compact-algebraic-diagrams}
is precisely the local compact-preservation condition. It places no
compactness restriction on the outer values $F(a)$.



The presheaf property of Lemma~\ref{lem:parameter-modules} proves
these assertions inductively. Local coefficient maps preserve compact
objects, giving the local compact-image conditions of
Definition~\ref{def:algebraic-diagrams}. Outer object values are
arbitrary in \eqref{eq:raw-presheaf-up}. Applying the construction in
coefficient parameter modules gives the full intertwining objects,
as in Remark~\ref{rem:parameter-detection}.


\begin{definition}\label{def:additive-relations}
Let $A\in\nPr{(n+1)}$. For each complete exact
sequence $e=(a\to b\to c)$ with compact vertices
in $A$, define a map $\theta_e$ in $\mathsf P_n(A)$
recursively.

For $n=1$, put
\[
 \theta_e:
 y_b\amalg_{y_a}y_0\longrightarrow y_c.
\]

For $n=s+1\geq2$, let $e'$ be the central exact
sequence in $\PrHom_A(b,b)$, and write
\[
 \theta_{e'}:x_{e'}\longrightarrow z_{e'}
 \quad\text{in }\mathcal R_A(b,b)
 =\mathsf P_s(\PrHom_A(b,b)).
\]
Put
\[
 J_s=\Fun([1]^{\op},\nPr{(s-1)}),
 \qquad I_s=\one_s,
\]
and let $\varepsilon_s:J_s\to I_s$ be evaluation
at $0$. Since the endpoints of $\theta_{e'}$
are compact, it determines a compact-preserving
coefficient map
\(
 J_s\longrightarrow\mathcal R_A(b,b)\).
By enriched Yoneda and tensor--Hom adjunction,
this corresponds to
\(
 \widetilde\theta_{e'}:
 J_s\odot y_b\longrightarrow y_b\).
Define
\[
 \theta_e:y_b\longrightarrow
 z_e:=y_b\amalg_{J_s\odot y_b}(I_s\odot y_b),
\]
where the other map in the pushout is
$\varepsilon_s\odot y_b$.

For each $m\geq0$, choose a small set
$\mathcal G_m$ representing the compact objects
of $\nPr m$, including the unit. Define
\begin{equation}\label{eq:additive-relation-set}
\begin{split}
 S_n(A)={}&
 \{g\odot(y_0\to0)\mid g\in\mathcal G_{n-1}\}\\
 &{}\cup
 \{g\odot\theta_e\mid
 e\in\EqRow_{n+1}^\omega(A),\
 g\in\mathcal G_{n-1}\}.
\end{split}
\end{equation}
\end{definition}

Choosing representatives of the exact sequences makes the indexing
collections small.


\begin{definition}\label{def:additive-motive}
Define
\begin{equation}\label{eq:presheaf-motive-definition}
 \M_n(A)=\mathsf P_n(A)[S_n(A)^{-1}],
 \qquad L_A:\mathsf P_n(A)\longrightarrow\M_n(A).
\end{equation}
We call $\M_n(A)$ the additive motive of $A$.
\end{definition}

The universality is proved in Theorem \ref{thm:compact-presheaf-up}.


\begin{lemma}\label{lem:compact-module-localization}
Every map in $S_n(A)$ has compact source and target. The localization
\eqref{eq:presheaf-motive-definition} is coefficient-linear, belongs to
$\nPr n$, and preserves compact objects. The objects $g\odot L_Ay_a$,
for $a\in A^\omega$ and compact coefficient $g$, compactly generate its
underlying presentable category.
\end{lemma}


\begin{proof}
The maps $\theta_e$ have compact endpoints
by induction: the base case uses finite colimits
of representables, while the higher cases use
compactness of $J_s$, $I_s$, and their actions
on representables.
The strongly saturated class generated by
$S_n(A)$ is stable under compact coefficient
actions, hence under all coefficient actions
by filtered colimits. Thus the localization
is coefficient-linear.

Since the relations have compact endpoints,
local objects are closed under filtered colimits,
so $L_A$ preserves compact objects.
The objects $g\odot L_Ay_a$, for compact $g$,
are compact generators, proving
$\M_n(A)\in\nPr n$.
\end{proof}


\subsection{The universal property}

\begin{lemma}\label{lem:outgoing-coinverter}
Let $E\in\nPr{(s+1)}$, and let an arrow $\alpha$ of
$\PrHom_E(x,y)$ be classified by $J_s\odot x\to y$. Then
\[
 y\longrightarrow y\amalg_{J_s\odot x}(I_s\odot x)
\]
is an equivalence if and only if $\alpha$ is an equivalence.
\end{lemma}

\begin{proof}
Taking coefficient Hom into $t\in E$ turns the morphism into
\begin{align}\label{equation:evalutaion-map-at-t}
 \PrHom_E(y,t)
 \times_{\PrHom_E(x,t)^{[1]}}\PrHom_E(x,t)\to \PrHom_E(y,t).
\end{align}
The map from the second factor \(\PrHom_E(x,t)\) of the pullback to the arrow category \(\PrHom_E(x,t)^{[1]}\) is the constant-arrow inclusion. Its
image is the equivalence arrows, and each fiber of \(\PrHom_E(x,t)\to \PrHom_E(x,t)^{[1]}\) is either empty or contractible. If $\alpha$ is invertible, every postcomposition lies in that image, so the pullback projection is an equivalence. Testing
against all coefficients and using ordinary Yoneda proves that the
original pushout map is an equivalence. Conversely, assume that the morphism in the statement is an equivalence.
Take $t=y$ and
lift the identity of $y$ along the map in
\eqref{equation:evalutaion-map-at-t}. Its image
in the arrow parameter identifies $\alpha$ with an equivalence arrow.
\end{proof}

\begin{theorem}\label{thm:compact-presheaf-up}\label{thm:presheaf-universal}
Restriction to $a\mapsto L_Ay_a$ gives a natural equivalence
\begin{equation}\label{eq:compact-presheaf-up}
 \Fun^L_{\nPr{(n-1)}}(\M_n(A),E)
 \simeq \Add_n^\omega(A,E)\cong \Add_{n,\mathrm{filt}}(A,E).
\end{equation}
It is compatible with every coefficient-enriched parameter.
\end{theorem}

\begin{proof}
The second equivalence is due to Proposition~\ref{prop:global-additive-extension}.
Let us prove the first equivalence.
By \eqref{eq:raw-presheaf-up}, a coefficient-linear
left adjoint $T:\mathsf P_n(A)\to E$ corresponds to
a diagram $F\in\operatorname{Fun}_n^\omega(A,E)$.
By coefficient-linearity, $T$ inverts the defining
relations \(S_n(A)\) precisely when $F(0)=0$ and $T(\theta_e)$
is an equivalence for every compact exact sequence $e$.

By the universality of localization, it suffices to prove that, under the condition \(F(0_A)=0_E\), \(T(\theta_e)\) is an equivalence if and only if \(F(e)\) is exact.
We prove by induction on $n$.
For $n=1$, this is the cofiber comparison.
For $n=s+1\geq2$, let $e=(a\to b\to c)$ and let
$e'$ be its central exact sequence in $\PrHom_A(b,b)$.
The construction of $\theta_e$ and
Lemma~\ref{lem:outgoing-coinverter} give that \(
 T(\theta_e)\) is an equivalence if and only if \(\Phi_{bb}(\theta_{e'})\) is an equivalence, where
\[
 \Phi_{bb}:\mathsf P_s(\PrHom_A(b,b))
 \longrightarrow\PrHom_E(Fb,Fb)
\]
is the induced coefficient-linear functor.
Its associated diagram preserves zero, since
zero morphisms factor through $0_A$ and $F(0_A)=0_E$.
By induction, the latter condition is equivalent to
exactness of the image of $e'$.
The recursive functoriality of $F$ preserves
adjunctions and their coherent data, so the
image of $e'$ is the central sequence of $F(e)$.
Thus its exactness is equivalent to that of $F(e)$.
\end{proof}







\begin{lemma}\label{lem:module-row-preservation}
The construction $\M_n$ takes internal exact sequences of coefficient
modules and compact-preserving module maps to internal exact sequences
of output modules.
\end{lemma}

\begin{proof}
Parameter functoriality transports the complete duality diagrams,
their units, counits, and mates. On a represented generator $U_H(h)$
of the middle output, the central exact sequence is the image under $U_H$ of
the source central exact sequence evaluated at $h$. It is exact by
Theorem~\ref{thm:presheaf-universal}. The same holds after every
compact coefficient action. All functors involved preserve colimits,
so the base cofiber identities extend from these generators to all
objects. Induction through the transported central diagrams proves
recursive exactness. The proof in every parameter module gives the
same assertion for complete coefficient diagrams.
\end{proof}




\begin{lemma}\label{lem:motive-filtered-colimits}
The functor $\M_n$ preserves filtered colimits
of coefficient modules and compact-preserving
module functors.
\end{lemma}

\begin{proof}
Compact objects and their coefficient Homs
descend along filtered colimits by
Lemma~\ref{lem:compact-coefficient-homs}.
The same holds for compact exact sequences
by Proposition~\ref{prop:fnd-classifier}.
Consequently,
\[
 \Add_n^\omega(\operatorname*{colim}_i A_i,E)
 \simeq\lim_i\Add_n^\omega(A_i,E).
\]
The same argument applies after every coefficient
parameter. Restricting to diagrams with compact
outer values gives the corresponding equivalence
for compact-preserving functors, since compact
objects descend along filtered colimits.
The assertion follows from the universal property
of colimits in the coefficient category.
\end{proof}



\subsection{Tensor products and functoriality}

\begin{proposition}\label{prop:motivic-products}
Compact-preserving coefficient-module maps induce functors on the
$\M_n$, naturally with respect to coefficient parameter diagrams.
There are external products
\[
 \M_n(A)\otimes\M_n(B)\longrightarrow\M_n(A\otimes B)
\]
preserving compact objects, with coherent associativity, symmetry,
and unit. 
In particular, for any $k\geq1$ and any $\mathbb E_k$-algebra
object $A$, Day convolution
endows $\M_n(A)$ with an $\mathbb E_k$-algebra
structure with compact unit $U_A(\one_A)$,
and $U_A$ is $\mathbb E_k$-monoidal.
\end{proposition}


\begin{proof}
The raw external products of
Definition~\ref{def:raw-presheaves}
make $\mathsf P_n$ lax symmetric monoidal.
For compact $b\in B$, the functor
$-\otimes b:A\to A\otimes B$
preserves compact objects and exact sequences.
Thus the composite
\[
 \mathsf P_n(A)\otimes\mathsf P_n(B)
 \longrightarrow\mathsf P_n(A\otimes B)
 \longrightarrow\M_n(A\otimes B)
\]
inverts the defining relations in the first
variable when the second is represented.
The same holds in the other variable.
Since representables and their coefficient actions
generate, the product descends to
\[
 \M_n(A)\otimes\M_n(B)\longrightarrow\M_n(A\otimes B).
\]
It preserves compact generators, and its
coherences follow from those of the raw product
and the parameterized universal property.
The last assertion is a formal fact of lax symmetric monoidal functors.
\end{proof}



\subsection{Comparison with relative BGT motive}

Let $a$ be a small idempotent-complete symmetric
monoidal stable category whose tensor product is
exact in each variable. An $a$-linear stable category
is an object of $\Mod_a(\Cat^{\mathrm{perf}})$.
Its small stable category of linear exact functors
is denoted by
$\Fun_a^{\ex}(\mathcal C,\mathcal D)$.
The base $a$ is called rigid if every object of $a$
is dualizable.


Assume that $a$ is rigid. The category
$\Mod_a(\Cat^{\mathrm{perf}})$ is the category
$\Cat^{\mathrm{perf}}(a)$ of
\citep[Definition~4.3]{HSS}.
Passing to Ind-completions makes it an object of
$\nPr2$, with coefficient Homs
\begin{equation}\label{eq:linear-Hom}
 \PrHom_{\Mod_a(\Cat^{\mathrm{perf}})}
       (\mathcal C,\mathcal D)
 \simeq\Ind\Fun_a^{\ex}(\mathcal C,\mathcal D).
\end{equation}
Indeed, free $a$-modules on compact stable
categories are compact generators, and the
coefficient action preserves compact objects.

\begin{definition}\label{def:relative-additive-motives}
Following \citep[Definition~5.3]{HSS}, a sequence
of $a$-linear stable categories is split-exact
if its underlying sequence is split-exact in
the sense of \citep[Definition~5.18]{BGT}.
Write $\operatorname{Mot}_{\mathrm{add}}^a$
for the category of additive $a$-motives
defined in \citep[Definition~5.8]{HSS}.
Its universal property is given by
\citep[Theorem~5.12]{HSS}.
\end{definition}

\begin{theorem}\label{thm:linear-normalization}
For every rigid $a$ as above, there is a
canonical equivalence
\begin{equation}\label{eq:BGT-normalization}
 \M_1(\Mod_a(\Cat^{\mathrm{perf}}))
 \simeq\operatorname{Mot}_{\mathrm{add}}^a.
\end{equation}
\end{theorem}


\begin{proof}
Since $a$ is rigid, right adjoints of
$a$-linear exact functors are canonically
$a$-linear, by applying
\citep[Proposition~3.11]{BBJ}
after Ind-completion.
By \citep[Proposition~5.4]{HSS}, relative
split-exact sequences are precisely the
localization sequences in the associated
linear $(\infty,2)$-category.
The adjunctions and the central cofiber
sequence $ir\to1\to jp$ identify these
with the exact sequences of
Definition~\ref{def:genuine}.
The equivalence follows from the universal
properties of
Theorem~\ref{thm:presheaf-universal}
and \citep[Theorem~5.12]{HSS}.
\end{proof}


For $a=\Sp^\omega$, this recovers the BGT
additive motive category \citep{BGT}.
The rigidity hypothesis is essential:
without it, an underlying split-exact
sequence need not be internally exact.